\topic{lumped-dacs}{Humble's lumped-parameter ballistics, worked through, and the model for a solid divert and attitude control system}

Reading anchors: Humble, Section~6.5.1, ``Lumped-Parameter Methods'',
printed pp.~334--341 / PDF pp.~354--361, Eqs.~(6.26)--(6.42)
\cite{humble}; Humble, Section~3.2.2, ``First Law for a Control
Volume'', printed pp.~83--86, Eqs.~(3.11)--(3.25); Hill and Peterson,
Section~12.7, Eq.~(12.28), printed p.~599 \cite{hillpeterson};
Williams, Section~9.4, Eqs.~(82)--(86), printed pp.~339--341
\cite{williams}; Sutton and Biblarz, Section~12.5, printed p.~483,
and Section~14.3, printed p.~541 \cite{sutton}.

Humble's Section~6.5.1 has two halves. The steady half (printed
pp.~334--337) sets the storage term to zero. The unsteady half
(``Unsteady Lumped-Parameter Method'', from printed p.~337) keeps it.
A divert and attitude control system (DACS) moves valves on the
time scale on which the chamber fills and empties, so the unsteady half
is the one we need. We work through both, because the steady half
supplies the equilibrium that the unsteady solutions approach and a
worked example we can check exactly.

This topic reuses the lumped chamber of
Section~\ref{topic:solid-ballistics} and its notation
(Section~\ref{sec:sb-notation}). It adds four things: the full algebra
of Humble's equations with his notation translated at each step, closed
forms for arbitrary valve schedules, the energy equation that Humble's
isothermal model omits, and the resulting DACS model.

\subsection{Humble's notation against ours}
\label{sec:ld-notation}
Humble writes Chapter~6 in the symbols of the first column. The second
column is our symbol for the same quantity. Parallel symbols of the
other books are listed in Section~\ref{sec:sb-mass-notation}.
\begin{center}
\begin{tabular}{p{2.2cm}p{2.0cm}p{7.6cm}}
Humble & Ours & Meaning, units, and differences\\ \hline
$p_c$ & $p_c$ & chamber pressure, Pa; Humble quotes MPa in
                 the burn law\\
$r_b$ & $r_b$ & burning (recession) rate; Humble's example uses cm/s\\
$a$, $n$ & $a$, $n$ & burn-law coefficient and exponent,
                 $r_b=ap_c^n$ (Humble Eq.~6.18); Humble's $a$ is in
                 cm/s/MPa$^n$\\
$\rho_p$ & $\rho_p$ & solid propellant density\\
$A_b$, $A_t$ & $A_b$, $A_t$ & burning area; throat area\\
$c^*$ & $c^*$ & characteristic velocity\\
$m$ & $m_c$ & gas mass stored in the chamber, kg\\
$V$ (also $V_c$) & $V_c$ & gas volume of the chamber, m$^3$\\
$\rho_g$ & $\rho_g$ & chamber gas density, $m/V$\\
$\dot m_{\mathrm{in}}$ & $\dot m_g$ & gas generated at the burning
                 surface, $r_b\rho_pA_b$\\
$\dot m_{\mathrm{out}}$ & $\dot m_{\mathrm{out}}$ & gas leaving
                 through the throat(s)\\
$T_c$ & $T_c$ & chamber gas temperature, K; the subscript $c$
                 means chamber\\
$T_{\mathrm{out}}$ & $T_{\mathrm{out}}$ & temperature of the gas
                 leaving the control volume\\
$c_p$ & $c_p$ & specific heat at constant pressure, J/(kg\,K)\\
$\mathcal M$ & $10^3N_0M$ & molar mass in kg/kmol; $M$ is the mass of
                 one molecule (Kittel--Kroemer)\\
$\mathcal R_u$ & $10^3R$ & universal constant, 8314 J/(kmol\,K);
                 our $R=N_0k_B$ is per mole\\
$\mathcal R_u/\mathcal M$ & $R_s$ & specific gas constant,
                 $R_s=k_B/M=R/(N_0M)$, J/(kg\,K); Humble's Chapter~3
                 calls this $R$ (his Eq.~3.27)\\
$\mathcal W$ & $y$ & web distance burned, normal to the surface\\
$\tau$ & $t_R$ & Humble's ``motor time constant''
                 $\rho_gV/\dot m_i$, a residence time (Turns' $t_R$);
                 our $\tau$ is $k_BT$, and our relaxation time is
                 $t_{\mathrm{rel}}$, Eq.~\eqref{eq:sb-tau}\\
$p'$, $k_b$, $k_t$ & $p'$, $k_b$, $k_t$ & $p_c/p_{c,i}$,
                 $A_b/A_{b,i}$, $A_t/A_{t,i}$ (dimensionless)\\
subscript $i$ & subscript $i$ & Humble's reference (``initial'')
                 state; in the DACS section $i$ also indexes valves,
                 and we then write the reference state with
                 subscript $\mathrm{ref}$
\end{tabular}
\end{center}
Our constants and thermodynamic conventions are those of
Section~\ref{sec:sb-notation}, from Kittel and Kroemer
\cite{CKittelHKroemer1980}: $R$ is the molar gas constant $N_0k_B$ and
$R_s=k_B/M$ is the specific one. Humble's $\mathcal R_u$ and
$\mathcal M$ appear only where his equations are quoted.
Two points need care. First, $T_c$ is not a static temperature: the
gas in the lumped chamber is at rest, so static and stagnation values
coincide there, and its outlets carry $T_{\mathrm{out}}=T_c$.
Second, Humble's $\tau$ is a
residence time, our $t_R$; we reserve $\tau=k_BT$ for temperature in
energy units. Our $t_{\mathrm{rel}}$ in Eq.~\eqref{eq:sb-tau} is the relaxation
time of the linearized pressure; Section~\ref{sec:ld-dimensionless}
shows $t_{\mathrm{rel}}=t_R/(1-n)$. From here on each Humble equation is
displayed first in his symbols and then in ours.

\subsection{The steady lumped method, printed pp.~334--337}
\label{sec:ld-steady}
\paragraph{Mass balance.} Humble's control volume (his Fig.~6.25)
is the gas in the chamber, fed by every exposed burning surface. His
Eqs.~(6.26)--(6.27) are
\begin{align*}
 \dot m_{\mathrm{out}}&=\frac{p_cA_t}{c^*},\\
 \frac{dm}{dt}&=\dot m_{\mathrm{in}}-\dot m_{\mathrm{out}}
            =r_b\rho_pA_b-p_c\frac{A_t}{c^*}.
\end{align*}
In our symbols, with $m\mapsto m_c$ and
$\dot m_{\mathrm{in}}\mapsto\dot m_g$,
\begin{equation}\label{eq:ld-mass}
 \frac{dm_c}{dt}=\dot m_g-\dot m_{\mathrm{out}}
 =\rho_pA_b r_b-\frac{p_cA_t}{c^*}.
\end{equation}
The outlet term is the choked-nozzle law, Eq.~\eqref{eq:sb-nozzle-flow}.
Equation~\eqref{eq:ld-mass} is exact under the lumped (B1) and
choked-outlet (B3) assumptions of Section~\ref{sec:sb-pressure-evolution};
it needs no assumption about temperature.

\paragraph{Steady state, Eqs.~(6.28)--(6.29).}
Humble's steady assumption is $dm/dt=0$. He states its meaning:
``the chamber pressure adjusts instantaneously to any changes in
burning rate, burning surface area, or other parameters.''
Setting the left side of \eqref{eq:ld-mass} to zero gives his
Eq.~(6.28),
\[
 r_b\rho_pA_b=p_c\frac{A_t}{c^*}
 \qquad\Longleftrightarrow\qquad
 \rho_pA_b r_b=\frac{p_cA_t}{c^*}.
\]
Insert the burn law $r_b=ap_c^n$ (Humble Eq.~6.18, our
Eq.~\eqref{eq:sb-burn-law}):
\[
 \rho_pA_b\,a\,p_c^n=\frac{p_cA_t}{c^*}.
\]
Multiply by $c^*/A_t$ and divide by $p_c^n>0$:
\[
 p_c^{1-n}=\frac{a\rho_pA_bc^*}{A_t}.
\]
For $n\neq1$ raise both sides to the power $1/(1-n)$. This is Humble's
Eq.~(6.29),
\begin{equation}\label{eq:ld-psteady}
 p_c=\left(\frac{a\rho_pA_bc^*}{A_t}\right)^{1/(1-n)},
\end{equation}
identical to our Eq.~\eqref{eq:sb-peq} with $K=A_b/A_t$.
Proposition~\ref{prop:sb-equilibrium} proves that it is the only
positive equilibrium when $A_b$ and $A_t$ are held fixed.

\paragraph{Units in Eq.~(6.29).} Humble gives $a$ in
cm/s/MPa$^n$, so $r_b$ in cm/s when $p_c$ is in MPa. Keep that
convention and write $p_c=\hat p\ \mathrm{MPa}$ with $\hat p$ a pure
number. The balance $\rho_pA_ba\hat p^{\,n}=\hat p\,(10^6\
\mathrm{Pa})A_t/c^*$ then reads
\[
 \hat p^{\,1-n}
 =\frac{\rho_pA_bc^*\,a}{10^6\ \mathrm{Pa}\,A_t},
\]
with $a$ in m/s. Every unit on the right cancels: kg\,m$^{-3}$ times
m/s times m/s over Pa is kg\,m$^{-1}$\,s$^{-2}$ over Pa, which is 1.
This is Humble's instruction that the group
$\rho_pA_bc^*/A_t$ be expressed in MPa/(cm/s).

\paragraph{Humble's worked example, Eq.~(6.30).} A cylindrical port
(his Fig.~6.21) has inner radius $r_i=10$~cm, outer radius
$r_o=20$~cm, length $L=200$~cm, throat radius 5~cm,
$c^*=1500$~m/s, $\rho_p=0.0015$~kg/cm$^3=1500$~kg/m$^3$, and
$r_b=0.5\,\hat p^{\,n}$~cm/s, i.e.\ $a=0.005$~m/s. Only the port surface
burns, so after a web distance $\mathcal W$ (our $y$) its area is
\[
 A_b=2\pi L\,(r_i+\mathcal W),\qquad
 A_t=\pi(5\ \mathrm{cm})^2=25\pi\ \mathrm{cm}^2.
\]
Their ratio, with $\mathcal W$ in cm, is
\[
 \frac{A_b}{A_t}
 =\frac{2\pi(200)(10+\mathcal W)}{25\pi}
 =16\,(10+\mathcal W).
\]
Substitute into the unit-consistent form above:
\begin{align*}
 \hat p^{\,1-n}
 &=\frac{(1500)(1500)(0.005)}{10^6}\cdot16\,(10+\mathcal W)\\
 &=0.01125\cdot16\,(10+\mathcal W)\\
 &=0.18\,(10+\mathcal W).
\end{align*}
This is Humble's Eq.~(6.30),
$p_c=[0.18(10+\mathcal W)]^{1/(1-n)}$, and it confirms his coefficient.

\paragraph{The web history, exactly.} Humble integrates
$\mathcal W=\int_0^{t}r_b\,dt$ (his Eq.~6.31) by the Euler steps of
his Fig.~6.26. The quasi-steady equation can be integrated in closed
form. With $\mathcal W$ in cm and $t$ in s,
\[
 \frac{d\mathcal W}{dt}=0.5\,\hat p^{\,n}
 =0.5\,[0.18\,(10+\mathcal W)]^{q},
 \qquad q\equiv\frac{n}{1-n}.
\]
Separate the variables:
\[
 (10+\mathcal W)^{-q}\,d\mathcal W=0.5\,(0.18)^q\,dt .
\]
For $q\neq1$ (that is, $n\neq\tfrac12$), integrate from $0$ to
$\mathcal W$ on the left and from $0$ to $t$ on the right:
\[
 \frac{(10+\mathcal W)^{1-q}-10^{1-q}}{1-q}=0.5\,(0.18)^q\,t .
\]
For $n=\tfrac12$ ($q=1$) the left side is
$\ln[(10+\mathcal W)/10]$ instead. The burn ends at $\mathcal W=w$,
the web thickness. The stated radii give $w=r_o-r_i=10$~cm, so
\begin{align*}
 t_b&=\frac{20^{1-q}-10^{1-q}}{0.5\,(1-q)(0.18)^q}
 \quad(n\neq\tfrac12),\\
 t_b&=\frac{\ln2}{0.5\cdot0.18}=7.70\ \mathrm s
 \quad(n=\tfrac12).
\end{align*}
The same exponents give the end pressures
$[0.18\cdot10]^{1/(1-n)}$ and $[0.18\cdot20]^{1/(1-n)}$:
\begin{center}
\begin{tabular}{cccc}
$n$ & $t_b$ (s) & $p_c(0)$ (MPa) & $p_c(t_b^-)$ (MPa)\\ \hline
0.3 & 13.22 & 2.32 & 6.23\\
0.4 & 10.54 & 2.66 & 8.46\\
0.5 & 7.70 & 3.24 & 12.96
\end{tabular}
\end{center}
These agree with the curves of Humble's Fig.~6.27 (burnout near
13.2, 10.5 and 7.7~s). The text on printed p.~336 says the total web is
40~cm. That is inconsistent with $r_o-r_i=10$~cm and with his own
figure; a 40~cm web would give $t_b\approx41$~s for $n=0.3$. We use
10~cm.

\paragraph{Why the pressure integral does not depend on $n$.}
Humble notes under Fig.~6.27 that the area under each curve is the
same. Integrate \eqref{eq:ld-mass} with $dm_c/dt=0$ over the burn:
the throat passes the whole propellant mass $m_p$,
\[
 \int_0^{t_b}\frac{p_cA_t}{c^*}\,dt=m_p
 \quad\Longrightarrow\quad
 \int_0^{t_b}p_c\,dt=\frac{c^*m_p}{A_t},
\]
which contains no $n$. In the example
$m_p=\rho_p\pi L(r_o^2-r_i^2)=282.7$~kg, so
$\int p_c\,dt=1500\cdot282.7/(7.854\times10^{-3})\ \mathrm{Pa\,s}
=54.0$~MPa\,s. The same number follows from the web variable,
using $dt=d\mathcal W/r_b$ and $\hat p/r_b=\hat p^{\,1-n}/0.5$:
\[
 \int\hat p\,dt=\int_0^{10}\frac{0.18(10+\mathcal W)}{0.5}\,d\mathcal W
 =0.36\left[10\mathcal W+\frac{\mathcal W^2}{2}\right]_0^{10}=54.0 .
\]

\paragraph{Energy in the steady method, Eq.~(6.32).}
Humble asks how the steady problem is solved without an energy
equation. For an adiabatic, calorically perfect gas he writes
$\dot m_{\mathrm{in}}c_pT_c-\dot m_{\mathrm{out}}c_pT_{\mathrm{out}}=0$;
with $\dot m_{\mathrm{in}}=\dot m_{\mathrm{out}}$ this gives
$T_{\mathrm{out}}=T_c$: the gas leaves at the chamber stagnation
temperature. Our symbols are the same here. Section~\ref{sec:ld-energy}
derives the unsteady energy equation and shows when $T_c$ may be held
at its steady value during transients.

\subsection{The unsteady lumped method, printed pp.~337--341}
\label{sec:ld-unsteady}
\paragraph{Accumulation, Eq.~(6.33).} Humble writes the stored mass
from the perfect-gas law,
\[
 m=\frac{p_cV\mathcal M}{\mathcal R_uT_c}
 \qquad\Longleftrightarrow\qquad
 m_c=\frac{p_cV_c}{R_sT_c}.
\]
Take logarithms, $\ln m_c=\ln p_c+\ln V_c-\ln R_s-\ln T_c$, and
differentiate in time. If all four factors may vary,
\begin{equation}\label{eq:ld-logdiff-full}
 \frac1{m_c}\frac{dm_c}{dt}
 =\frac1{p_c}\frac{dp_c}{dt}+\frac1{V_c}\frac{dV_c}{dt}
  -\frac1{R_s}\frac{dR_s}{dt}-\frac1{T_c}\frac{dT_c}{dt}.
\end{equation}
Humble's Eq.~(6.34) keeps only the first two terms on the right:
\[
 \frac1m\frac{dm}{dt}=\frac1{p_c}\frac{dp_c}{dt}+\frac1V\frac{dV}{dt}.
\]
Dropping the last two terms of \eqref{eq:ld-logdiff-full} is his
assumption that $\mathcal M$ and $T_c$ stay fixed while $p_c$ changes.
In our notation, fixed $\mathcal M$ is fixed molecular mass $M$, hence
fixed $R_s=k_B/M$.
He cites Heister and Landsbaum \cite{heisterlandsbaum1991}: these
variations may be neglected ``unless a very rapid and large pressure
disturbance is present,'' because ``the energy of adiabatic compression
\ldots is balanced by an increase in energy flux through the throat.''
Section~\ref{sec:ld-energy} tests this statement.

\paragraph{Volume growth, Eq.~(6.35), and the pressure equation,
Eq.~(6.36).} The receding surface uncovers volume at the rate
$dV/dt=r_bA_b$, our Eq.~\eqref{eq:sb-volume-rate} times $dy/dt=r_b$.
Solve Humble's (6.34) for the pressure derivative and substitute
(6.27) and (6.35), one term at a time:
\begin{align*}
 \frac1{p_c}\frac{dp_c}{dt}
 &=\frac1m\frac{dm}{dt}-\frac1V\frac{dV}{dt}\\
 &=\frac1m\left(r_b\rho_pA_b-\frac{p_cA_t}{c^*}\right)
   -\frac{r_bA_b}{V}\\
 &=\frac1m\left(r_b\rho_pA_b-\frac{p_cA_t}{c^*}
   -\frac mV\,r_bA_b\right)\\
 &=\frac1m\left[r_bA_b(\rho_p-\rho_g)-p_c\frac{A_t}{c^*}\right].
\end{align*}
The third line factors out $1/m$; the fourth uses $\rho_g=m/V$.
The last line is Humble's Eq.~(6.36). In our symbols, and multiplying
by $p_c/m_c=R_sT_c/V_c$,
\begin{equation}\label{eq:ld-pressure}
 \frac{dp_c}{dt}
 =\frac{R_sT_c}{V_c}
  \left[(\rho_p-\rho_g)A_b\,a\,p_c^n-\frac{p_cA_t}{c^*}\right],
\end{equation}
which is our Eq.~\eqref{eq:sb-pressure-ode} and Hill and Peterson's
Eq.~(12.28). Humble adds that $\rho_p/\rho_g\approx300$, so the
$\rho_g$ correction is usually dropped. Section~\ref{topic:solid-ballistics}
estimates it at about 0.6\% of the equilibrium pressure.

\subsection{Dimensionless form and the motor time constant}
\label{sec:ld-dimensionless}
\paragraph{Deriving Eq.~(6.37).} Humble states the dimensionless form
and refers to Heister and Landsbaum for the derivation; here it is.
Drop $\rho_g$ in \eqref{eq:ld-pressure}, keep $V_c$ and $T_c$ fixed over
the interval considered, and let $A_b(t)$ and $A_t(t)$ vary. Choose a
reference equilibrium state, Humble's subscript $i$, where
\begin{equation}\label{eq:ld-ref-flow}
 \dot m_i=a\,p_{c,i}^{\,n}A_{b,i}\rho_p=\frac{p_{c,i}A_{t,i}}{c^*}.
\end{equation}
Define $p'=p_c/p_{c,i}$, $k_b=A_b/A_{b,i}$ and $k_t=A_t/A_{t,i}$.
Substitute $p_c=p_{c,i}p'$, $A_b=A_{b,i}k_b$ and $A_t=A_{t,i}k_t$
into the bracket of \eqref{eq:ld-pressure}:
\begin{align*}
 \rho_pA_bap_c^n-\frac{p_cA_t}{c^*}
 &=\rho_pA_{b,i}k_b\,a\,p_{c,i}^{\,n}p'^{\,n}
   -\frac{p_{c,i}A_{t,i}}{c^*}k_tp'\\
 &=\dot m_i\,k_bp'^{\,n}-\dot m_i\,k_tp'
   \qquad\text{by \eqref{eq:ld-ref-flow}}.
\end{align*}
The left side of \eqref{eq:ld-pressure} is
$dp_c/dt=p_{c,i}\,dp'/dt$. Divide by $p_{c,i}$:
\[
 \frac{dp'}{dt}=\frac{R_sT_c\,\dot m_i}{V_cp_{c,i}}
   \left(k_bp'^{\,n}-k_tp'\right).
\]
The reference gas density is $\rho_{g,i}=p_{c,i}/(R_sT_c)$, so the
prefactor is $\dot m_i/(\rho_{g,i}V_c)$. Define
\begin{equation}\label{eq:ld-tauH}
 t_R\equiv\frac{\rho_{g,i}V_c}{\dot m_i}
 =\frac{m_{c,i}}{\dot m_i}.
\end{equation}
Then
\begin{equation}\label{eq:ld-dimensionless}
 \boxed{t_R\frac{dp'}{dt}=k_b(t)\,p'^{\,n}-k_t(t)\,p'.}
\end{equation}
This is Humble's Eq.~(6.37). We have written $k_b(t)$ and $k_t(t)$
as functions of time: nothing in the derivation required them to be
constant. For a DACS, $k_t(t)$ is the valve command.

\paragraph{What $t_R$ is.} Equation~\eqref{eq:ld-tauH} is the time
to empty the reference gas inventory at the reference flow: a
residence time (Turns calls the same ratio $t_R=\rho V/\dot m$, his
Eq.~(6.37), p.~196 \cite{turns}). Substitute
$\dot m_i=p_{c,i}A_{t,i}/c^*$ and $R_sT_c=(\Gamma c^*)^2$ from
Eq.~\eqref{eq:sb-nozzle-flow}:
\[
 t_R=\frac{p_{c,i}}{R_sT_c}\,V_c\,\frac{c^*}{p_{c,i}A_{t,i}}
 =\frac{V_c\,c^*}{R_sT_cA_{t,i}}
 =\frac{L^*}{\Gamma^2c^*},
 \qquad L^*=\frac{V_c}{A_{t,i}} .
\]
The pressure cancels, so $t_R$ depends only on $L^*$ and the gas.
In energy units the gas enters through $(\Gamma c^*)^2=R_sT_c=\tau_c/M$,
with $\tau_c=k_BT_c$.
Our relaxation time, Eq.~\eqref{eq:sb-tau}, is
$t_{\mathrm{rel}}=L^*/[(1-n)\Gamma^2c^*]$, hence
\begin{equation}\label{eq:ld-tau-relation}
 t_{\mathrm{rel}}=\frac{t_R}{1-n}.
\end{equation}
Humble quotes $0.01<t_R<0.5$~s for most solid motors.

\subsection{Closed-form solutions of Eq.~(6.37)}
\label{sec:ld-closed}
\paragraph{The substitution that linearizes it.} Equation
\eqref{eq:ld-dimensionless} is a Bernoulli equation. Let
\[
 u\equiv p'^{\,1-n}>0 .
\]
By the chain rule,
$du/dt=(1-n)\,p'^{\,-n}\,dp'/dt$. Multiply
\eqref{eq:ld-dimensionless} by $(1-n)p'^{\,-n}/t_R$:
\begin{align}
 \frac{du}{dt}
 &=\frac{1-n}{t_R}\,p'^{\,-n}\left(k_bp'^{\,n}-k_tp'\right)
 \notag\\
 &=\frac{1-n}{t_R}\left(k_b-k_t\,p'^{\,1-n}\right)
 \notag\\
 &=\frac{1-n}{t_R}\left(k_b(t)-k_t(t)\,u\right).
 \label{eq:ld-linear}
\end{align}
This is \emph{linear} in $u$, for any time histories $k_b(t)$ and
$k_t(t)$. Every solution below follows from it.

\paragraph{Constant $k_b$ and $k_t$: Eq.~(6.40).} With constant
coefficients, \eqref{eq:ld-linear} has the equilibrium $u_\infty=k_b/k_t$,
and $u-u_\infty$ obeys
$d(u-u_\infty)/dt=-(1-n)k_t(u-u_\infty)/t_R$. With $u(0)=1$
(the reference state),
\[
 u(t)=\frac{k_b}{k_t}
      +\left(1-\frac{k_b}{k_t}\right)e^{-(1-n)k_tt/t_R}
     =\frac{k_b}{k_t}-\frac{k_b-k_t}{k_t}\,e^{-(1-n)k_tt/t_R}.
\]
Undo the substitution, $p'=u^{1/(1-n)}$:
\begin{equation}\label{eq:ld-640}
 p'(t)=\left[\frac{k_b}{k_t}
       -\frac{k_b-k_t}{k_t}\,e^{-(1-n)k_tt/t_R}\right]^{1/(1-n)},
\end{equation}
Humble's Eq.~(6.40). It approaches
$p'_\infty=(k_b/k_t)^{1/(1-n)}$, which is
\eqref{eq:ld-psteady} at the new areas divided by its reference value.

\paragraph{A step in burning area: Eq.~(6.41).} Put $k_t=1$:
\[
 p'=\left[k_b-(k_b-1)\,e^{-(1-n)t/t_R}\right]^{1/(1-n)}.
\]
The effective time constant is $t_R/(1-n)=t_{\mathrm{rel}}$, as in
Proposition~\ref{prop:sb-stability}.

\paragraph{A step in throat area: Eq.~(6.42), corrected.} Put
$k_b=1$ in \eqref{eq:ld-640} and factor $1/k_t$ out of the bracket:
\begin{align}
 p'&=\left[\frac1{k_t}-\frac{1-k_t}{k_t}\,e^{-(1-n)k_tt/t_R}
     \right]^{1/(1-n)}\notag\\
   &=\left(\frac1{k_t}\right)^{1/(1-n)}
     \left[1-(1-k_t)\,e^{-(1-n)k_tt/t_R}\right]^{1/(1-n)}.
 \label{eq:ld-642}
\end{align}
Humble prints Eq.~(6.42) without the exponent $1/(1-n)$ on the second
bracket, and introduces it with ``Eq.~(6.41) gives'' where Eq.~(6.40)
is meant (we checked the page image, PDF p.~361). The printed form
fails at $t=0$: it gives $k_t^{-1/(1-n)}\cdot k_t=k_t^{-n/(1-n)}$
instead of $p'(0)=1$. For $k_t=0.75$ and $n=0.35$ that is 1.168
instead of 1. Equation~\eqref{eq:ld-642} satisfies $p'(0)=1$ and agrees
with direct numerical integration of \eqref{eq:ld-dimensionless} to
$3\times10^{-12}$ (a finite check with $t_R=0.02$~s and
$k_t\in\{0.5,0.75,1.1\}$). Humble's text and final value
$k_t^{-1/(1-n)}$ are correct.

The effective time constant is $t_R/[k_t(1-n)]$: it depends on the
\emph{new} throat area. Closing valves ($k_t<1$) makes the pressure
settle more slowly; opening them makes it settle faster.

\paragraph{Blowdown, Eqs.~(6.38)--(6.39).} After burnout
$k_b=0$. With $k_t=1$, Eq.~\eqref{eq:ld-dimensionless} is
$t_R\,dp'/dt=-p'$, so
\[
 \frac{dp'}{p'}=-\frac{dt}{t_R},\qquad
 \frac{p_c}{p_{c,i}}=e^{-t/t_R}.
\]
There is no factor $1-n$: without generation the burn law plays no
part. Humble calls this ``isothermal blowdown'' and says that for an
adiabatic, reversible blowdown one should use $T_c\propto
p^{(\gamma-1)/\gamma}$ in Eq.~(6.37). We carry that out. With no
generation and fixed $V_c$, $dm_c/dt=V_c\,d\rho_g/dt=-\dot m_{\mathrm{out}}$.
Isentropic expansion gives
$\rho_g=\rho_{g,i}\,p'^{\,1/\gamma}$ and
$T_c=T_{c,i}\,p'^{\,(\gamma-1)/\gamma}$. The choked flow is
$\dot m_{\mathrm{out}}=\Gamma p_cA_t/\sqrt{R_sT}$, so
\[
 \dot m_{\mathrm{out}}
 =\frac{\Gamma p_{c,i}A_t}{\sqrt{R_sT_{c,i}}}\,
  p'\,p'^{\,-(\gamma-1)/(2\gamma)}
 =\dot m_i\,p'^{\,(\gamma+1)/(2\gamma)} .
\]
The storage term is
$V_c\,d\rho_g/dt=V_c\rho_{g,i}\,\gamma^{-1}p'^{\,(1-\gamma)/\gamma}\,dp'/dt$.
Equate the two and divide by $\dot m_i$, using
$t_R=\rho_{g,i}V_c/\dot m_i$:
\[
 \frac{t_R}{\gamma}\,p'^{\,(1-\gamma)/\gamma}\frac{dp'}{dt}
 =-p'^{\,(\gamma+1)/(2\gamma)}
 \quad\Longrightarrow\quad
 \frac{dp'}{dt}=-\frac{\gamma}{t_R}\,p'^{\,(3\gamma-1)/(2\gamma)} .
\]
The exponents combine as
$(\gamma+1)/(2\gamma)-(1-\gamma)/\gamma=(3\gamma-1)/(2\gamma)$.
Separate variables; since $1-(3\gamma-1)/(2\gamma)=-(\gamma-1)/(2\gamma)$,
\[
 p'^{\,-(\gamma-1)/(2\gamma)}-1=\frac{\gamma-1}{2\gamma}\cdot
 \frac{\gamma t}{t_R},
\]
so
\begin{equation}\label{eq:ld-adiabatic-blowdown}
 p'(t)=\left[1+\frac{(\gamma-1)\,t}{2t_R}\right]^{-2\gamma/(\gamma-1)}.
\end{equation}
For small $t$, $p'\approx1-\gamma t/t_R$: the adiabatic pressure
first falls $\gamma$ times faster than the isothermal one, because
the gas also cools. With $\gamma=1.2$ at $t=t_R$ the two give 0.319
(adiabatic) and 0.368 (isothermal). Equation
\eqref{eq:ld-adiabatic-blowdown} agrees with direct integration of the
mass and energy equations of Section~\ref{sec:ld-energy} to
$2\times10^{-12}$.

\paragraph{Arbitrary valve schedules.} Return to the linear equation
\eqref{eq:ld-linear} with $k_t(t)$ an arbitrary piecewise-continuous
command. Define
\[
 \Phi(t)\equiv\frac{1-n}{t_R}\int_0^tk_t(s)\,ds .
\]
Multiply \eqref{eq:ld-linear} by $e^{\Phi(t)}$. The left side becomes
$d(e^{\Phi}u)/dt$, because $d\Phi/dt=(1-n)k_t/t_R$:
\[
 \frac{d}{dt}\left(e^{\Phi}u\right)
 =e^{\Phi}\frac{du}{dt}+e^{\Phi}\frac{1-n}{t_R}k_tu
 =e^{\Phi}\frac{1-n}{t_R}k_b .
\]
Integrate from 0 to $t$ with $u(0)=u_0$:
\begin{equation}\label{eq:ld-integrating-factor}
 u(t)=e^{-\Phi(t)}\left[u_0+\frac{1-n}{t_R}
      \int_0^tk_b(s)\,e^{\Phi(s)}\,ds\right],
 \qquad p'=u^{1/(1-n)} .
\end{equation}
For piecewise-constant commands, as when valves are pulsed, apply
\eqref{eq:ld-640} on each interval with the value of $u$ at the end of
the previous interval as its starting value:
\[
 u(t_{j}+s)=\frac{k_b}{k_{t,j}}
 +\left(u(t_j)-\frac{k_b}{k_{t,j}}\right)e^{-(1-n)k_{t,j}s/t_R},
 \quad 0\le s\le t_{j+1}-t_j .
\]
This is exact under the hypotheses of \eqref{eq:ld-dimensionless}.
For a square wave with $k_t=1$ for 4~ms and $0.6$ for 6~ms,
$t_R=0.02$~s and $n=0.35$, it agrees with a numerical solution to
$3\times10^{-8}$ over 20 periods. It is the golden vector for a pulsed
valve.

\paragraph{Fast pulsing: average and ripple.} Let $k_t$ alternate
between $k_1$ for a fraction $D$ of each period $T_p$ and $k_0$ for
the rest, with $k_b$ fixed. Suppose $T_p$ is short compared with
$t_R/[(1-n)k_t]$, so $u$ changes little within a period; call its
period average $\bar u$. In the periodic state, the rise of $u$ during
the $k_0$ phase equals the fall during the $k_1$ phase. By
\eqref{eq:ld-linear},
\[
 \frac{1-n}{t_R}(k_b-k_0\bar u)(1-D)T_p
 =\frac{1-n}{t_R}(k_1\bar u-k_b)DT_p .
\]
Collect the $\bar u$ terms:
$k_b(1-D)+k_bD=\bar u\,[Dk_1+(1-D)k_0]$, so
\[
 \bar u=\frac{k_b}{\bar k_t},\qquad
 \bar k_t\equiv Dk_1+(1-D)k_0 .
\]
The mean pressure is set by the time-averaged area,
$\bar p'\approx(k_b/\bar k_t)^{1/(1-n)}$. The peak-to-peak change of
$u$ is the rise during the $k_0$ phase. Substitute
$k_b=\bar k_t\bar u$ and $\bar k_t-k_0=D(k_1-k_0)$:
\[
 \delta u=\frac{1-n}{t_R}\,\bar u\,(\bar k_t-k_0)(1-D)T_p
 =\frac{1-n}{t_R}\,\bar u\,(k_1-k_0)\,D(1-D)\,T_p .
\]
Since $\delta p'/p'\approx\delta u/[(1-n)u]$, the relative pressure
ripple is
\begin{equation}\label{eq:ld-ripple}
 \frac{\delta p'}{\bar p'}\approx
 (k_1-k_0)\,D(1-D)\,\frac{T_p}{t_R},
\end{equation}
independent of $n$. For the square wave above,
$\bar p'=(1/0.76)^{1/0.65}=1.5253$ against a period mean of 1.5255
from the numerical solution, and \eqref{eq:ld-ripple} predicts a
peak-to-peak ripple of 4.8\% against 4.79\% computed.

\subsection{The energy equation Humble omits, and when isothermal is enough}
\label{sec:ld-energy}
\paragraph{Humble's first law for a control volume.} Humble's
Eq.~(3.25), printed p.~86, is
\[
 \dot Q=P_s+\frac{\partial}{\partial t}\int_{cv}e\rho\,dV
 +\int_{cs}\left(h+pe+\frac{v^2}{2}\right)\rho\,
   (\vec v\cdot\hat n)\,dA ,
\]
where $\dot Q$ is heat added, $P_s$ shaft and shear power, $e$ the
specific stored energy, $h=u+p\upsilon$ the enthalpy (his Eq.~3.24,
$\upsilon$ the specific volume), $pe$ the specific potential energy,
$\vec v$ the velocity and $\hat n$ the outward normal. His control
volume is fixed in space.

\paragraph{Choosing a fixed control volume.} Take the whole interior
of the case: gas plus remaining solid. It is fixed, as Eq.~(3.25)
requires, and its only open boundary is the throat. The solid is at
rest; its volume $V_s$ and the gas volume satisfy $V_s+V_c=$ const.
Neglect $P_s$ and potential energy. Assume a single gas of fixed
composition (fixed $M$) with constant $c_p$, $c_v$,
$R_s=c_p-c_v$ and $\gamma=c_p/c_v$ (Section~\ref{sec:sb-notation}).
Measure energies from the reference at which the gas at temperature
$T$ has $h=c_pT$ and $u=c_vT$, and give the solid the specific enthalpy
$h_s=c_pT_f$, where $T_f$ is the adiabatic flame temperature at
constant pressure: burning a unit of solid at constant pressure and
without heat loss produces a unit of gas at $T_f$. This is the
definition of $T_f$ (Humble Section~4.4; Turns Chapter~2), and $T_f$
is the $T_c$ of the steady model. The solid's specific internal
energy is $u_s=h_s-p_c/\rho_p$.

Evaluate each term of Eq.~(3.25). The gas stores $m_cc_vT_c$; the solid
stores $\rho_pV_su_s$, which changes only through $V_s$ because $u_s$
is fixed (we neglect the small $p_c$-dependence of $u_s$). Hence
\[
 \frac{\partial}{\partial t}\int_{cv}e\rho\,dV
 =\frac{d}{dt}(m_cc_vT_c)+\rho_pu_s\frac{dV_s}{dt}
 =\frac{d}{dt}(m_cc_vT_c)-\rho_pu_s\frac{dV_c}{dt}.
\]
At the throat the gas leaves with stagnation enthalpy
$h+v^2/2=c_pT_{\mathrm{out}}$, so the surface integral is
$c_pT_{\mathrm{out}}\dot m_{\mathrm{out}}$. With
$\rho_p\,dV_c/dt=\rho_pA_br_b=\dot m_g$,
\[
 \dot Q=\frac{d}{dt}(m_cc_vT_c)
        -\dot m_g\left(c_pT_f-\frac{p_c}{\rho_p}\right)
        +c_pT_{\mathrm{out}}\dot m_{\mathrm{out}} .
\]
Since $\dot m_gp_c/\rho_p=p_c\,dV_c/dt$, write
$E\equiv m_cc_vT_c$ and rearrange:
\begin{equation}\label{eq:ld-energy}
 \boxed{\frac{dE}{dt}=c_pT_f\,\dot m_g-p_c\frac{dV_c}{dt}
        -c_pT_{\mathrm{out}}\dot m_{\mathrm{out}}+\dot Q .}
\end{equation}
Generated gas enters with enthalpy $c_pT_f$; the term $-p_c\,dV_c/dt$
is the flow work spent filling the volume the grain uncovers. For the
lumped chamber, whose gas leaves at its own stagnation state,
$T_{\mathrm{out}}=T_c$. With \eqref{eq:ld-mass} and
$p_cV_c=m_cR_sT_c=(\gamma-1)E$, the state $(m_c,E)$ closes the model.
This is the chamber counterpart of the plenum equations already used in
\path{GasGeneratorValveSystem.h}.

\paragraph{Temperature and pressure forms.} Set
$T_{\mathrm{out}}=T_c$. Expand
$dE/dt=c_vT_c\,dm_c/dt+m_cc_v\,dT_c/dt$ and use \eqref{eq:ld-mass}:
\[
 m_cc_v\frac{dT_c}{dt}
 =\dot m_g(c_pT_f-c_vT_c)-p_c\frac{dV_c}{dt}
  -\dot m_{\mathrm{out}}(c_pT_c-c_vT_c)+\dot Q .
\]
Write $c_pT_f-c_vT_c=c_p(T_f-T_c)+R_sT_c$ and
$c_pT_c-c_vT_c=R_sT_c$, and use $p_c=\rho_gR_sT_c$:
\begin{align*}
 m_cc_v\frac{dT_c}{dt}
 &=\dot m_gc_p(T_f-T_c)\\
 &\quad+R_sT_c\left(\dot m_g-\dot m_{\mathrm{out}}
   -\rho_g\frac{dV_c}{dt}\right)+\dot Q\\
 &=\dot m_gc_p(T_f-T_c)+R_sT_cV_c\frac{d\rho_g}{dt}+\dot Q .
\end{align*}
The last line uses
$dm_c/dt=V_c\,d\rho_g/dt+\rho_g\,dV_c/dt$. So the temperature relaxes
toward $T_f$ through new gas, and rises when the gas is compressed
($d\rho_g/dt>0$). For the pressure, differentiate
$p_cV_c=(\gamma-1)E$ and substitute \eqref{eq:ld-energy}, using
$(\gamma-1)c_p=\gamma R_s$:
\begin{equation}\label{eq:ld-pressure-energy}
 V_c\frac{dp_c}{dt}
 =\gamma R_s\left(T_f\dot m_g-T_{\mathrm{out}}\dot m_{\mathrm{out}}\right)
  -\gamma p_c\frac{dV_c}{dt}+(\gamma-1)\dot Q .
\end{equation}
Compare the isothermal Eq.~\eqref{eq:sb-general-pressure-outlet}, in
which $T_c=T_{\mathrm{out}}=T_f$ throughout:
$V_c\,dp_c/dt=R_sT_f(\dot m_g-\dot m_{\mathrm{out}})-p_c\,dV_c/dt$.
The adiabatic form has an extra factor $\gamma$. Setting $T_c=T_f$ in
the temperature equation shows why: $T_c$ stays at $T_f$ without heat
exchange only if $d\rho_g/dt=0$, that is, at constant pressure. An
isothermal transient implicitly assumes heat exchange with the walls.

\paragraph{Linearization with the energy equation.} Freeze $V_c$,
$A_b$ and $A_t$, set $\dot Q=0$, neglect $dV_c/dt$, and write
$\dot m_{\mathrm{out}}=\Gamma p_cA_t/\sqrt{R_sT_{\mathrm{out}}}$ with
$T_{\mathrm{out}}=T_c$, so that the outflow depends on temperature. The equilibrium is $T_c=T_f$,
$\dot m_g=\dot m_{\mathrm{out}}=\dot m_i$. Perturb with
$P=\delta p_c/p_{c,i}$, $\Theta=\delta T_c/T_f$ and $s=t/t_R$.
To first order,
\[
 \frac{\delta\dot m_g}{\dot m_i}=nP,\qquad
 \frac{\delta\dot m_{\mathrm{out}}}{\dot m_i}=P-\frac{\Theta}{2}.
\]
Insert these into \eqref{eq:ld-pressure-energy}, whose equilibrium
part cancels, and use $\gamma R_sT_f\dot m_i/(V_cp_{c,i})=\gamma/t_R$:
\[
 \frac{dP}{ds}
 =\gamma\left[nP-\left(P-\frac\Theta2\right)-\Theta\right]
 =-\gamma(1-n)P-\frac{\gamma}{2}\Theta .
\]
The term $-\Theta$ comes from $-T_{\mathrm{out}}\dot m_{\mathrm{out}}$ with
$\delta T_c$. For the temperature, $\ln T_c=\ln p_c+\ln V_c-\ln R_s-\ln m_c$
gives $d\Theta/ds=dP/ds-\delta(dm_c/ds)/m_{c,i}$, and by
\eqref{eq:ld-mass}
$\delta(dm_c/ds)/m_{c,i}=nP-P+\Theta/2$. Therefore
\[
 \frac{d\Theta}{ds}
 =-\gamma(1-n)P-\frac\gamma2\Theta+(1-n)P-\frac\Theta2
 =-(\gamma-1)(1-n)P-\frac{\gamma+1}{2}\Theta .
\]
The system is $d(P,\Theta)^{\mathsf T}/ds=J(P,\Theta)^{\mathsf T}$ with
\[
 J=\begin{pmatrix}
 -\gamma(1-n) & -\gamma/2\\
 -(\gamma-1)(1-n) & -(\gamma+1)/2
 \end{pmatrix}.
\]
Its determinant and trace are
\begin{align*}
 \det J&=\frac{\gamma(1-n)(\gamma+1)}{2}-\frac{\gamma(\gamma-1)(1-n)}{2}
        =\gamma(1-n),\\
 \operatorname{tr}J&=-\gamma(1-n)-\frac{\gamma+1}{2}.
\end{align*}
The discriminant is
$\Delta=(\operatorname{tr}J)^2-4\det J
=\tfrac14\left[(3-2n)^2\gamma^2-2(5-6n)\gamma+1\right]$.

\begin{proposition}[energy-equation modes]\label{prop:ld-modes}
For $\gamma>1$ and $0\le n<1$ both eigenvalues of $J$ are real and
negative. Their product is $\gamma(1-n)$.
\end{proposition}
\begin{proof}
Negativity: $\det J>0$ and $\operatorname{tr}J<0$, so if the
eigenvalues are real both are negative, and if complex their real
part $\tfrac12\operatorname{tr}J$ is negative. Reality: $4\Delta$ is a
quadratic in $\gamma$ with positive leading coefficient $(3-2n)^2$. Its
own discriminant is
\[
 4(5-6n)^2-4(3-2n)^2=4\,(2-4n)(8-8n)=64\,(1-2n)(1-n).
\]
If $\tfrac12<n<1$ this is negative, so $4\Delta>0$ for every $\gamma$.
If $n=\tfrac12$ the only root is $\gamma=\tfrac12<1$. If
$0\le n<\tfrac12$ the larger root is
$\gamma_+=[(5-6n)+4s]/(3-2n)^2$ with $s=\sqrt{(1-2n)(1-n)}$. Then
$\gamma_+\le1$ is equivalent to $4s\le(3-2n)^2-(5-6n)=4-6n+4n^2$,
whose right side is positive (its discriminant $36-64$ is negative).
Squaring,
\[
 (4-6n+4n^2)^2-16(1-2n)(1-n)
 =4n^2(9-12n+4n^2)=4n^2(3-2n)^2\ge0 .
\]
So $\gamma_+\le1<\gamma$, and $4\Delta>0$.
\end{proof}
For $\gamma=1.2$ and $n=0.35$ the eigenvalues are $-0.618$ and $-1.262$
(in units of $1/t_R$); a finite-difference Jacobian of the full
$(m_c,E)$ equations reproduces them. Compare three models:
\begin{center}
\begin{tabular}{lll}
Model & Pressure decay rate $\times t_R$ & Source\\ \hline
isothermal & $1-n=0.65$ & Humble (6.37)\\
isentropic chamber, $\Theta$ ignored & $\gamma(1-n)=0.78$ &
  Williams (83)--(84)\\
mass and energy & $0.618$ and $1.262$ & \eqref{eq:ld-energy}
\end{tabular}
\end{center}
Williams writes $\omega_L=a^2C_D/L^*$ with $a^2=\gamma R_sT_c$ and
$C_D=\dot m_{\mathrm{out}}/(A_tp_c)=1/c^*$; his linearized Eq.~(83)
decays at $\omega_L(1-n)=\gamma(1-n)/t_R$ for a quasi-steady burn
law. The slow mode of the full model is within 5\% of Humble's
isothermal rate. The fast mode is temperature relaxation. This
quantifies Heister and Landsbaum's statement: the pressure transient
is nearly isothermal because hot new gas and the throat flow carry
away the compression energy, but a disturbance on the scale of
$t_R$ also excites a temperature mode. For a DACS, whose valves move
on that scale, carry $E$; it costs one extra state per volume.

\subsection{The model for a solid DACS}
\label{sec:ld-dacs}
\paragraph{Two architectures.} Sutton's Fig.~12--27 (printed p.~483)
feeds hot-gas valves directly from the grain's gas volume through
insulated manifolds; Fig.~12--28 does the same for an interceptor, with
divert nozzles through the centre of gravity and a separate attitude
set. Call this \emph{architecture~A}: one lumped volume, $N$ valves.
Huzel and Huang's gas generator (their p.~116) and our code put a
fixed orifice and a plenum between the grain and the valves:
\emph{architecture~B}. Both valve arrangements are normally open;
Sutton: ``the chamber pressure rises when any valve is closed.''

\paragraph{Valve flow and the valve command.} Valve $j$ has opening
$x_j(t)\in[0,1]$, geometric flow area $A_j(x_j)$ and discharge
coefficient $C_{D,j}$. (Our pintle law is
$A_j=\pi R_j^2[1-(1-x_j)^2]$, Section~6 of
\path{documents/derivations/SolidRocketMotorGasGenerator.md}; none of the six textbooks gives a pintle area law.) While
the valve is choked it passes $C_{D,j}A_jp_{\mathrm{up}}/c^*$, with
$p_{\mathrm{up}}$ the pressure of the volume it drains. For a vacuum
or high-altitude exit it stays choked. Define the effective total
\[
 A_v(t)\equiv\sum_{j=1}^{N}C_{D,j}A_j\bigl(x_j(t)\bigr).
\]
All choked valves on one volume share one upstream state, so their
outflows add to $A_v\,p_{\mathrm{up}}/c^*$ (Eq.~\eqref{eq:sb-general-outlets}).
Thrust and moment follow per valve:
\[
 \mathbf F=\sum_jF_j\,\hat{\mathbf e}_j,\qquad
 \mathbf M=\sum_j(\mathbf r_j-\mathbf r_{\mathrm{cg}})\times
           F_j\,\hat{\mathbf e}_j,\qquad
 F_j=C_{F,j}\,p_{\mathrm{up}}\,C_{D,j}A_j ,
\]
with $\hat{\mathbf e}_j$ the thrust direction and $C_{F,j}$ the
thrust coefficient at that valve's current area ratio and pressure
ratio.

\paragraph{Architecture A: Eq.~(6.37) with $k_t=A_v/A_{v,\mathrm{ref}}$.}
With the valves on the grain volume, Section~\ref{sec:ld-dimensionless}
applies directly with $A_t\mapsto A_v$. Four consequences follow.

\emph{Pressure amplification and valve cross-coupling.} Closing
valves raises the pressure to $k_t^{-1/(1-n)}$. With four equal valves
and $n=0.35$, closing one gives $k_t=3/4$ and $p'_\infty=1.557$;
closing two gives $k_t=1/2$ and $p'_\infty=2.905$. Every open valve's
thrust scales with $p_{\mathrm{up}}$, so each valve command changes
every other valve's thrust. The side force from closing valve 1 of an
opposed pair is the full thrust of valve 2 at the raised pressure.

\emph{Asymmetric response.} By \eqref{eq:ld-642} the settling time
after a step is $t_R/[k_t(1-n)]$ at the new $k_t$: closing is slower
than opening. For pulsed commands use the exact piecewise solution and
\eqref{eq:ld-ripple}.

\emph{An admissible band of total valve area.} The case limits
$p_c\le p_{\max}$. Sutton Section~14.3 gives a low-pressure
deflagration limit of 0.05--0.15~MPa, so we also need
$p_c\ge p_{\min}$. At equilibrium, \eqref{eq:ld-psteady} with
$A_t\mapsto A_v$ is decreasing in $A_v$ for $n<1$, so
\[
 p_{\min}\le p_{c,\infty}\le p_{\max}
 \iff
 a\rho_pA_bc^*\,p_{\max}^{\,n-1}\le A_v\le
 a\rho_pA_bc^*\,p_{\min}^{\,n-1}.
\]
The ratio of the bounds is $(p_{\max}/p_{\min})^{1-n}$. The
controller may not command total areas outside this band for longer
than the transient allows. All valves shut is outside it: $A_v=0$ has
no equilibrium, and the pressure grows until the case fails. Sutton
also notes (Section~14.3) that fast depressurization extinguishes a
propellant, with higher pressures needing faster $dp/dt$; this bounds
how fast valves may open. The book gives no threshold values.

\emph{$L^*$ instability at large opening.} $L^*=V_c/A_v$ falls as
valves open. Williams' stability limit (his Eq.~86, printed p.~341)
is a critical value of $L^*\bar p^{\,2n}$, below which the chamber
oscillates at low frequency and can turn into chuffing. Opening
valves lowers both $L^*$ and $\bar p$, so the fully open state is the
one to check. The critical value needs the burning-rate response of
Williams Section~9.1.5.4, which a quasi-steady $r_b=ap^n$ cannot supply.

\paragraph{Architecture B: grain, orifice, plenum, valves.} Each volume
carries mass and energy. With subscripts $c$ (grain chamber) and
$\mathrm{pl}$ (plenum), and $\dot m_{\mathrm{link}}$ the orifice flow,
\begin{align*}
 \frac{dm_c}{dt}&=\dot m_g-\dot m_{\mathrm{link}},\\
 \frac{dE_c}{dt}&=c_pT_f\dot m_g-p_c\frac{dV_c}{dt}
                -c_pT_{\mathrm{out},c}\,\dot m_{\mathrm{link}},\\
 \frac{dm_{\mathrm{pl}}}{dt}&=\dot m_{\mathrm{link}}-\sum_j\dot m_j,\\
 \frac{dE_{\mathrm{pl}}}{dt}&=c_pT_{\mathrm{out},c}\,\dot m_{\mathrm{link}}
   -c_pT_{\mathrm{out},\mathrm{pl}}\sum_j\dot m_j-\dot Q_{\mathrm{wall}} .
\end{align*}
The first and third are Eq.~\eqref{eq:sb-two-chamber-mass}; the
second is \eqref{eq:ld-energy}; the fourth applies Humble's Eq.~(3.25)
to the rigid plenum, whose inflow carries the grain chamber's
stagnation enthalpy (for reverse flow use the plenum's). Each volume's
outflow leaves at its own stagnation state, so
$T_{\mathrm{out},c}=T_c$ and $T_{\mathrm{out},\mathrm{pl}}=T_{\mathrm{pl}}$;
$E_c=m_cc_vT_c$ and $E_{\mathrm{pl}}=m_{\mathrm{pl}}c_vT_{\mathrm{pl}}$.

\emph{Choked orifice: the grain is decoupled.} While the orifice
(effective area $C_oA_o$) is choked,
$\dot m_{\mathrm{link}}=C_oA_op_c/c^*$ depends only on the grain
chamber, so valve motion does not reach the grain. The plenum then
obeys, isothermally at $T_{\mathrm{pl}}$,
\[
 \frac{V_{\mathrm{pl}}}{R_sT_{\mathrm{pl}}}\frac{dp_{\mathrm{pl}}}{dt}
 =\dot m_{\mathrm{link}}-\frac{A_v(t)\,p_{\mathrm{pl}}}{c^*},
\]
which is \emph{linear} in $p_{\mathrm{pl}}$: Eq.~\eqref{eq:ld-dimensionless}
with $n=0$. Its equilibrium is
$p_{\mathrm{pl},\infty}=c^*\dot m_{\mathrm{link}}/A_v$. In steady
state the plenum gas is at the grain temperature, so the two $c^*$
values agree and
\begin{equation}\label{eq:ld-plenum-ratio}
 \frac{p_{\mathrm{pl},\infty}}{p_c}=\frac{C_oA_o}{A_v}.
\end{equation}
Its time constant is $t_{R,\mathrm{pl}}=V_{\mathrm{pl}}c^*/(R_sT_{\mathrm{pl}}A_v)$
with no $1-n$ factor.

\emph{When the orifice unchokes.} The orifice stays choked while
$p_{\mathrm{pl}}/p_c\le\pi_{\mathrm{cr}}
=[2/(\gamma+1)]^{\gamma/(\gamma-1)}$, the isentropic throat pressure
ratio of Section~\ref{sec:sb-pressure-evolution}. With
\eqref{eq:ld-plenum-ratio}, it stays choked while
\[
 A_v\ \ge\ \frac{C_oA_o}{\pi_{\mathrm{cr}}}.
\]
Below that, the link flow depends on both pressures. For isentropic
flow from $(p_c,T_c)$ to a throat at $p_{\mathrm{pl}}$, the exit
temperature is $T_c\,\pi^{(\gamma-1)/\gamma}$ with
$\pi=p_{\mathrm{pl}}/p_c$, the speed is
$\sqrt{2c_pT_c(1-\pi^{(\gamma-1)/\gamma})}$ and the density is
$p_{\mathrm{pl}}/(R_sT_c\pi^{(\gamma-1)/\gamma})$. Their product gives
\[
 \dot m_{\mathrm{link}}
 =C_oA_o\,p_c\sqrt{\frac{2\gamma}{(\gamma-1)R_sT_c}
   \left(\pi^{2/\gamma}-\pi^{(\gamma+1)/\gamma}\right)} .
\]
The grain chamber then feels the valves again.

\emph{Check against our simulation.} In the scenario
\texttt{CartridgeWithFourPintleValves} ($\gamma=1.25$,
$\pi_{\mathrm{cr}}=0.555$, $C_oA_o=4.015\times10^{-4}$~m$^2$, four
valves of $C_{D}A=2.985\times10^{-4}$~m$^2$), Eq.~\eqref{eq:ld-plenum-ratio}
predicts $p_{\mathrm{pl}}/p_c=0.33631$ with four valves open. The
simulation gives $2.2962/6.8277=0.33631$ at $t=0.375$~s. The choking
bound is $A_v\ge7.23\times10^{-4}$~m$^2$: four open valves
($1.19\times10^{-3}$) satisfy it and two ($5.97\times10^{-4}$) do not.
That is why, after two valves close, the orifice unchokes and the
grain pressure, steady at 6.83~MPa before the closure, rises to
7.11~MPa by $t=0.575$~s in the simulation. Architecture~B
buys isolation of the grain from valve commands for exactly the range
$A_v\ge C_oA_o/\pi_{\mathrm{cr}}$, at the cost of the pressure drop
across the orifice.

\paragraph{What to model.} The state, inputs and outputs are:
\begin{itemize}
\item[D1] \emph{State}: web $y$; grain chamber $(m_c,E_c)$; for
architecture~B also plenum $(m_{\mathrm{pl}},E_{\mathrm{pl}})$;
optionally actuator states for $x_j$.
\item[D2] \emph{Geometry}: $A_b(y)$ and $V_c(y)$ with
$dV_c/dy=A_b$ (Eq.~\eqref{eq:sb-volume-rate}); $V_{\mathrm{pl}}$ fixed.
\item[D3] \emph{Generation}: $\dot m_g=\rho_pA_b\,a(T_b)\,p_c^n$,
entering with enthalpy $c_pT_f$.
\item[D4] \emph{Links}: orifice flow choked or unchoked as above; each
valve choked at $C_{D,j}A_j(x_j)p_{\mathrm{up}}/c^*$ with $c^*$ at the
upstream temperature.
\item[D5] \emph{Inputs}: valve commands $x_j(t)$, prescribed or from a
controller; grain temperature $T_b$; ambient pressure.
\item[D6] \emph{Outputs}: $p_c$, $p_{\mathrm{pl}}$, $T_c$,
$T_{\mathrm{pl}}$, per-valve thrust, $\mathbf F$ and $\mathbf M$,
expelled mass, and a mass residual.
\item[D7] \emph{Constraints to monitor}: $p_{\max}$, $p_{\min}$, the
area band above, the choking bound, the opening rate, and $L^*$ at
full opening.
\end{itemize}
The existing code has D1 without $E_c$ (the grain chamber is
isothermal at $T_c$), D2--D4 for tubular and end-burning grains, D5
as prescribed schedules, and a plenum energy equation. It lacks the
grain-chamber energy equation, valve thrust directions and moments,
constraint monitors, and closed-loop control.

\paragraph{Golden vectors from this topic.} Each is a finite check on
a specialization, not physical validation:
\begin{enumerate}
\item Humble's Eq.~(6.30) and the burn times in the table of
Section~\ref{sec:ld-steady}; $\int p_c\,dt=c^*m_p/A_t$.
\item Eqs.~\eqref{eq:ld-640} and \eqref{eq:ld-642} for steps in
$k_b$ and $k_t$ (single volume, fixed $V_c$, isothermal,
$\rho_g$ neglected).
\item The piecewise solution for pulsed valves, and
Eq.~\eqref{eq:ld-ripple} as a fast-pulsing limit.
\item Isothermal blowdown $e^{-t/t_R}$ and adiabatic blowdown
\eqref{eq:ld-adiabatic-blowdown} ($k_b=0$).
\item The eigenvalues of Proposition~\ref{prop:ld-modes} for the
linearized mass and energy equations.
\item Eq.~\eqref{eq:ld-plenum-ratio} and the choking bound for
architecture~B.
\end{enumerate}

\paragraph{What these books do not give.} A pintle area law and a
hot-gas valve discharge coefficient as functions of stroke; valve
actuator dynamics; a non-quasi-steady burning rate for fast pressure
changes beyond Williams' linear response (his Section~9.1.5.4);
extinction thresholds for $dp/dt$ (Sutton cites his Ref.~14--1); wall
heat loss in the plenum and manifolds; and two-phase flow through the
valves. Each needs outside data before the model can claim more than
the specializations above.
